\scanpage{198}
ALGÈBRES de LIE SEMI-SIMPLES \mspage{(6)}

Tribu n° 59

Un cor. à la prop. 1 dira : ou bien tous les $\mathrm{e}_{\mathrm n}$ sont lint. indép. (car de poids différents), ou bien l'un est nul et alors $\lambda$ est un entier $\mathrm{q}\geqslant 0$, indépendance jusqu'à $\mathrm{e}_{\mathrm q}$, nuls ensuite. Ceci remplace une partie de la prop. 2. Dans l'exer. du bas de la p. 26, ajouter la stabilité par les binomiaux $\binom{\mathrm{H}}{\mathrm{n}}$.

\uline{n° 2} - Titre : les modules simples $\mathrm{W}(\mathrm{m})$, (dire ``simple'' au lieu de ``irréductible''). Dans la prop. 3 convenir que $\mathrm{e}_{-1}=0$ et $\mathrm{e}_{\mathrm n}=0$ pour $\mathrm{n}\geqslant\mathrm{m}+1$. Prendre $\mathrm{e}=\alpha^{\mathrm m}/\mathrm{m}!$ ; alors $\mathrm{e}_{\mathrm n}=\mathrm{Y}^{\mathrm n}\mathrm{e}/\mathrm{n}!$ ; qualifier la base $(\mathrm{e}_{\mathrm n})$ de ``canonique'' ; arranger (4).

Remplacer l'exer. 2, (p. 28), par un truc de Gabriel. Mettre l'exer. 3, (p. 29), dans le texte. Remonter ici le lemme 4 des pp. 53-54.

\uline{n° 3} - Titre : décomposition des $\underline{\mathrm{s}}$-modules. Ajouter ``de dim. finie'' dans la prop. 5 et son cor. Il semble qu'on n'a pas besoin du (ii) et du (iii) de la prop. 6 ; si oui, ça passe en exer. et on bloque le (i) avec la prop. 5. Dire ``dim. finie'' dans la prop. 7.

A la p. 30, l.12, donner une référence pour la complète réductibilité. Mieux détailler la remarque 1), p. 31. A la place de la remarque 2), (pas si ``facile''), on donnera le calcul du Casimir, (cf. Papa Cartan).

\uline{n° 4} - Dire, (l.5, p. 31), qu'on a $\mathrm{x}^2=0$. Dire que $\mathrm{E}$ est de dim. finie dans la prop. 8.On préfère la seconde démonstration p. 34, (les exponentielles de loct. nilpotents seront définies en Alg. IV). A la fin de la remarque, (p. 34 bas), dire : ceci montre aussi que ce sont les seules représentations irréd. holomorphes de $\mathrm{G}$. Expliciter la fin de la démonstration de la prop. 9. Lire `` $\mathrm{E}=\mathrm{S}^{\mathrm m}(\mathrm{V})$ '' à la l.4 bas de la p. 35. Dans la prop. 11, a) dire que $\unsure{\vartheta|\theta}^4=1$ ;

% archive sheet number: 197
